\originalchapter{动量、冲量与质心}\label{ch:momentum}
\originalsection{动量与冲量}
用同样的力推空车和满车, 速度变化不同. 对恒质量物体, 动量定义为 $\vct p=m\vct v$, 则牛顿第二定律可写为 $\dot{\vct p}=\vct F$. 动量是向量, 其单位 kg m/s. 冲量 $\vct J=\int_{t_1}^{t_2}\vct F\dd t$ 累积力的时间作用, 从运动方程积分得到 $\Delta\vct p=\vct J$. 它适用于力随时间变化、甚至碰撞过程短到不能细描加速度的情形.

碰撞中忽略重力冲量要求碰撞时间短且接触冲量远大于 $mg\Delta t$, 不是说重力不存在. 弹性球以速率 $u$ 垂直落地并以同速反弹, 取向上为正, 动量变化为 $2mu$; 粘在地面则为 $mu$. 平均接触力还需除以碰撞持续时间. 不同过程有相同动量变化也可能有不同峰值力.

\originalsection{系统与质心}
对包含固定一组质点的封闭系统, 总动量 $\vct P=\sum_i m_i\vct v_i$. 成对内力抵消, 得 $\dot{\vct P}=\vct F_{\rm ext}$. 外力为零时总动量守恒, 即使每个质点之间有强烈作用. 若只在某个方向上外力为零, 仅能保证该方向动量守恒.

\begin{definition}
总质量 $M=\sum m_i$ 的质心位置为 $\vct R=M^{-1}\sum_i m_i\vct r_i$. 连续分布对应 $\vct R=M^{-1}\int\vct r\dd m$, 其中 $\dd m$ 是所选微小质量元.
\end{definition}
恒定各质量时求导得到 $M\dot{\vct R}=\vct P$, 再求导得到 $M\ddot{\vct R}=\vct F_{\rm ext}$. 这说明质心按总外力运动; 并不说明各物体都围着质心静止. 爆炸抛出碎片时质心仍沿原抛物线, 前提是空气作用可忽略.

均匀细杆长度 $L$ 沿 $x$ 轴放置, $\dd m=(M/L)\dd x$, 所以 $X=M^{-1}\int_0^Lx(M/L)\dd x=L/2$. 非均匀密度 $\lambda(x)$ 则分母和分子都要积分. 有空洞的板可先算完整板质量矩再减去被挖部分, “负质量”在此只是积分差的记账, 不是真实负质量.

\originalsection{质心系与相对动能}
取 $\vct r_i=\vct R+\vct r'_i$, 有 $\sum m_i\vct r'_i=0$ 和 $\sum m_i\vct v'_i=0$. 展开动能,
\begin{equation}
K=\tfrac12\sum m_i|\vct V+\vct v'_i|^2
=\tfrac12M V^2+\tfrac12\sum m_i v_i'^2.\label{eq:cmenergy}
\end{equation}
混合项因质心定义消失. 第一项为质心平动能, 第二项为相对质心的内部运动能. 两物体质量 $m_1,m_2$ 时, 设相对速度 $\vct u=\vct v_1-\vct v_2$, 可由质心速度公式求出 $\vct v'_1=(m_2/M)\vct u$, $\vct v'_2=-(m_1/M)\vct u$, 因而内部动能为 $\mu u^2/2$, $\mu=m_1m_2/M$ 称为约化质量.

另一惯性系相对原系匀速 $\vct w$ 时, $K'=K-\vct w\cdot\vct P+Mw^2/2$. 单个状态动能依赖参考系, 但孤立碰撞前后 $M$、$\vct P$ 不变, 所以 $\Delta K'=\Delta K$. 判断碰撞是否弹性不会因改换惯性系而改变.
